% ---------------------------------------------------------------
% Preamble for "Hardware from the Ground Up"
% Compile with:  tectonic hardware-from-the-ground-up.tex
% ---------------------------------------------------------------
\usepackage[letterpaper,margin=1in,headheight=14pt]{geometry}
\usepackage{amsmath,amssymb,mathtools}
\usepackage{siunitx}
\usepackage[american]{circuitikz}
\usetikzlibrary{arrows.meta,positioning,calc,shapes.geometric,decorations.pathreplacing,patterns,automata}
\usepackage{tikz-timing}
\usepackage{booktabs,tabularx,array,multirow}
\usepackage{enumitem}
\usepackage{listings}
\usepackage{xcolor}
\usepackage{tcolorbox}
\tcbuselibrary{skins,breakable}
\usepackage{fancyhdr}
\usepackage{microtype}
\usepackage[hidelinks]{hyperref}
\usepackage{xurl}
\setlength{\emergencystretch}{2em}

\raggedbottom
\setlength{\parskip}{0.45em}
\setlength{\parindent}{0pt}
\renewcommand{\arraystretch}{1.25}
\setlist{itemsep=0.15em,topsep=0.3em}

% ----- headers --------------------------------------------------
\pagestyle{fancy}
\fancyhf{}
\fancyhead[L]{\small\nouppercase{\leftmark}}
\fancyhead[R]{\small\thepage}
\renewcommand{\headrulewidth}{0.4pt}
\fancypagestyle{plain}{\fancyhf{}\fancyfoot[C]{\small\thepage}\renewcommand{\headrulewidth}{0pt}}
\makeatletter
\renewcommand{\chaptermark}[1]{\markboth{\ifnum\value{chapter}>0 \@chapapp\ \thechapter.\ \fi #1}{}}
\makeatother

% wider number boxes in the table of contents (sections like 20.13)
\makeatletter
\renewcommand*\l@section{\@dottedtocline{1}{1.5em}{2.9em}}
\renewcommand*\l@subsection{\@dottedtocline{2}{4.4em}{3.6em}}
\makeatother

% ----- handy macros ----------------------------------------------
\newcommand{\VDD}{V_{DD}}
\newcommand{\Vt}{V_{t}}
\newcommand{\Vtn}{V_{tn}}
\newcommand{\Vtp}{V_{tp}}
\newcommand{\nt}[1]{\overline{#1}}
\newcommand{\code}[1]{\texttt{#1}}
\newcommand{\term}[1]{\textbf{#1}}

% ----- boxes -----------------------------------------------------
\newtcolorbox{bigidea}{enhanced,breakable,colback=blue!3,colframe=blue!45!black,
  fonttitle=\bfseries,title={The big idea},left=6pt,right=6pt}
\newtcolorbox{everyday}{enhanced,breakable,colback=teal!3,colframe=teal!50!black,
  fonttitle=\bfseries,title={An everyday picture},left=6pt,right=6pt}
\newtcolorbox{trap}{enhanced,breakable,colback=red!3,colframe=red!55!black,
  fonttitle=\bfseries,title={Common trap},left=6pt,right=6pt}
\newtcolorbox{onoa}{enhanced,breakable,colback=orange!6,colframe=orange!75!black,
  fonttitle=\bfseries,title={Seen on the assessment},left=6pt,right=6pt}
\newtcolorbox{worked}[1]{enhanced,breakable,colback=green!3,colframe=green!40!black,
  fonttitle=\bfseries,title={Worked example: #1},left=6pt,right=6pt}
\newtcolorbox{defn}[1]{enhanced,breakable,colback=gray!4,colframe=gray!55!black,
  fonttitle=\bfseries,title={#1},left=6pt,right=6pt}
\newtcolorbox{remember}{enhanced,breakable,colback=yellow!6,colframe=yellow!45!black,
  fonttitle=\bfseries,title={Remember},left=6pt,right=6pt}

% ----- code listings ----------------------------------------------
\lstdefinestyle{book}{
  basicstyle=\ttfamily\small,
  keywordstyle=\color{blue!60!black}\bfseries,
  commentstyle=\color{green!40!black}\itshape,
  stringstyle=\color{red!60!black},
  numbers=left, numberstyle=\tiny\color{gray},
  frame=single, framerule=0.4pt, rulecolor=\color{gray!50},
  backgroundcolor=\color{gray!5},
  breaklines=true, columns=fullflexible, keepspaces=true,
  showstringspaces=false, xleftmargin=1.8em, framexleftmargin=1.6em,
  aboveskip=0.8em, belowskip=0.8em}
\lstset{style=book}

% ----- reusable drawings ------------------------------------------
% 2:1 multiplexer.  #1 = name, #2 = position, #3 = label inside.
% Coordinates created: #1-in0, #1-in1, #1-sel, #1-out
\newcommand{\muxtwo}[3]{%
  \begin{scope}[shift={#2}]
    \draw[thick] (0,0) -- (0,2) -- (1,1.6) -- (1,0.4) -- cycle;
    \node[font=\scriptsize] at (0.18,1.5) {0};
    \node[font=\scriptsize] at (0.18,0.5) {1};
    \node[font=\small] at (0.6,1) {#3};
    \coordinate (#1-in0) at (0,1.5);
    \coordinate (#1-in1) at (0,0.5);
    \coordinate (#1-sel) at (0.5,0.2);
    \coordinate (#1-out) at (1,1);
  \end{scope}}

% D flip-flop / latch box. #1 = name, #2 = position, #3 = title above box,
% #4 = 1 to draw the edge triangle (flip-flop), 0 for a latch (level input).
% Coordinates: #1-d, #1-q, #1-clk
\newcommand{\dbox}[4]{%
  \begin{scope}[shift={#2}]
    \draw[thick] (0,0) rectangle (1.5,2);
    \node[anchor=west,font=\small] at (0.05,1.5) {D};
    \node[anchor=east,font=\small] at (1.45,1.5) {Q};
    \ifnum#4=1
      \draw (0,0.32) -- (0.22,0.5) -- (0,0.68);
      \node[anchor=west,font=\scriptsize] at (0.24,0.5) {CLK};
    \else
      \node[anchor=west,font=\scriptsize] at (0.05,0.5) {EN};
    \fi
    \node[font=\small\bfseries] at (0.75,2.3) {#3};
    \coordinate (#1-d) at (0,1.5);
    \coordinate (#1-q) at (1.5,1.5);
    \coordinate (#1-clk) at (0,0.5);
  \end{scope}}

% Timing diagrams: pass [xscale=..., yscale=...] per diagram.
